% 30_dworkin 部：Michelman 410--419（Michelman 第 IV 部 A--C 節）。
% 頁碼依 PDF 頁減 396：411（PDF 15）、412（16）、413（17）、414（18）、415（19）、416（20）、417（21）、418（22）、419（23）；註 53 跨 415--416，兩頁圖已核。
% Dworkin 本人著作（Freedom's Law、Life's Dominion、A Matter of Principle、Law's Empire）與 Rawls、Post《Constitutional Domains》本機無藏本，一律轉引 Michelman 註文。
% 判決之 U.S. Reports 卷頁取自 Michelman 註 3、37、38、47。
\chapter{Dworkin 之實質民主}
\label{ch:30_dworkin}

\begin{qaquestion}
Michelman\ix{persons}{michelman@Michelman, Frank I.} 在原文第 IV 部 A 至 C 節（411--418 頁）轉述並檢討 Dworkin\ix{persons}{dworkin@Dworkin, Ronald} 對 Brennan\ix{persons}{brennan@Brennan, William J., Jr.} 與民主\ix{terms}{01@\zhlidxcat{民主理論}!民主}之間關係的答案。
\begin{enumerate}
\tightlist
\item Dworkin 以憲法的內容判斷其民主性，這個論證對 accountability 民主觀與自治\ix{terms}{01@\zhlidxcat{民主理論}!自治}民主觀各生何種後果？
\item 道德解讀\ix{terms}{08@\zhlidxcat{憲法解釋與司法自制}!道德解讀}以及合理多元\ix{terms}{02@\zhlidxcat{正當性與合理歧見}!合理多元}的事實，如何使 Michelman 排除以制憲與修憲程序證立憲法之民主性？
\item 註 53 所列三種反對各主張什麼，Michelman 如何回應？
\item Dworkin 所列成員身分三項條件，Michelman 為何認為不足以使個人成為自己的立法者？
\end{enumerate}
\end{qaquestion}

\section{以內容判斷憲法之民主性}
\label{sec:30_dworkin-1}

Michelman 411 頁轉述 Dworkin 的論證：在根本法\ix{terms}{03@\zhlidxcat{憲政主義與憲法之效力基礎}!根本法}（laws of lawmaking，規範造法之法\ix{terms}{03@\zhlidxcat{憲政主義與憲法之效力基礎}!造法之法}）的層次，民主與否取決於這些法律實際說了什麼，無從以誰制定來回答。Brown v. Board of Education\ix{cases}{02@美國!01@聯邦最高法院!1954Brown v. Board of Education（1954）@Brown v. Board of Education（1954）}（347 U.S. 483, 1954）的內容有利於民主，Plessy v. Ferguson\ix{cases}{02@美國!01@聯邦最高法院!1896Plessy v. Ferguson（1896）@Plessy v. Ferguson（1896）}（163 U.S. 537, 1896）的內容則否。根本法的內容有利於民主之目的與價值，須具備下列各項：排除種姓（caste）、保障廣泛而公平的政治參與權（franchise）、防止恣意的法律差別待遇與其他壓迫性的國家權力運用、確保政府尊重思想、表達與結社的自由以及每位公民的智識與道德獨立。\footcite[411 \& n.39]{michelman1998}

Michelman 把這個主張歸納為：憲法可以因內容而具民主性，一如它可以因作者而具民主性。他評為「unassailable as far as it goes」〔就其所及而言無可攻擊〕。其推論如下：一旦承認民主的目標在使根本法的選擇合乎對民主政體的正確或最佳構想，並承認確有這樣一種正確或最佳構想（Dworkin 承認，註 41），那麼獨立司法機關所達成的憲法解釋\ix{terms}{08@\zhlidxcat{憲法解釋與司法自制}!解釋}，實際上就可能比人民與其代言人所提出的更接近該構想；讓法官參與治理因此不必然反民主，讓與法院最終的憲法解釋權也不能「in principle」〔就原則而言〕歸為反民主。\footcite[411--412 \& nn.41--42]{michelman1998}Michelman 補充，Dworkin 所辯護的對象，是 Brennan 的生涯所代表的裁判方式，Brennan 其人不在其內。\footcite[412]{michelman1998}

註 40 指出 Ely\ix{persons}{ely@Ely, John Hart} 在 \en{Democracy and Distrust}（1980）等著作中提出過類似論證，並附帶說明 Dworkin 與 Ely 的憲法理論有「important differences」〔重要差異〕，指引讀者參看 Dworkin \en{A Matter of Principle} 33--71 頁。差異的內容，Michelman 在該註並未展開。\footcite[411 n.40]{michelman1998}

\section{Accountability 民主觀與自治民主觀}
\label{sec:30_dworkin-2}

Michelman 412 頁分兩種聽眾。認為民主的用處只在 accountability 與福祉的人，會輕易接受 Dworkin 的論證：美國的分權政府包括由司法機關部分治理，自始即以較能回應被治理者的真實利益為由受到辯護，註 44 列舉 Madison 在 \en{The Federalist} 第 10、46、51 號與 Hamilton 在第 78 號的論述。Dworkin 把同一思路延伸到人民對本國民主政治狀態的利益，Michelman 在註 43 說明這項利益包括人民參與政治言論、結社、投票與擔任公職的樂趣。分權政府，連同司法機關的部分治理，可能比較純粹的程序民主更能實現這項利益。\footcite[412 \& nn.43--44]{michelman1998}

對 Dworkin 本人，Michelman 認為事情不止於此，因為在他那裡民主之要旨除 accountability 與福祉之外尚有自治。自治是活動，也是程序，採自治觀的民主，似乎免不了以程序性的檢驗判斷民主是否存在。自治觀下的民主，在一國人民以民主的政治程序自行決定一切政治上可決定的生活條件時最為圓滿。\footcite[412]{michelman1998}這與僅以內容判斷民主性的主張如何並存，Michelman 說並不清楚。

自治觀是否要求人民決定一切政治上可決定的問題，是否須延伸到決定根本法，亦即設定政府權力的目標與界限、並建立其後一切立法之制度的規範，Michelman 在 413 頁的答案是「inescapably」〔無從迴避〕的肯定，理由是人民有道德上與物質上的理由重視這些法律的內容。他隨即指出，人們總會有一種衝動，要把根本法劃出\ix{terms}{03@\zhlidxcat{憲政主義與憲法之效力基礎}!劃出}民主程序的管轄範圍：民主程序只決定根本法所留下的其餘政治選擇，根本法本身交給另一個「pre-democratic authority」〔先於民主的權威〕，由它運用道德理性，不訴諸政治意志，決定正義、繁榮、自由、乃至民主本身的根本法前提。Dworkin 主張以內容判斷憲法的民主資格，不問其作者，在 Michelman 看來正是這種衝動的一例。衝動背後有一道看似壓倒性的邏輯反對，指向把根本法的內容交由民主程序決定的構想，Michelman 留待原文第 V 部處理，第 \ref{ch:40_paradox} 部（涵蓋原文第 V、VI 部）依其論證展開。\footcite[413]{michelman1998}

\section{道德解讀與程序性證立}
\label{sec:30_dworkin-3}

Michelman 413 頁另舉一個理由，說明 Dworkin 何以在憲法的內容而不在其制定程序中尋找民主資格。有些論者主張憲法的民主性以程序證立：人民透過批准與修憲的政治程序，直接選擇了根本法。Michelman 以為這條路對 Dworkin 以及一切與他同樣堅持法律具解釋性格的人封閉。理由在於，憲法在適用之時，必是某人對構成它的文本、傳統與先例所作解釋的產物。Dworkin 將此論點推到不少人不願跟進的地步，主張對憲法中抽象的權利宣示用語，如言論自由、自由、正當程序、平等保護，不能不作「moral reading」〔道德解讀\ix{terms}{08@\zhlidxcat{憲法解釋與司法自制}!道德解讀}〕：解釋者須把這些用語當作政治道德價值或原則的「invocations」〔援引〕，從美國憲法傳統中視為主要固定點之處提煉出來，而提煉免不了引入解釋者自己對憲政之目的與理想的實質見解。\footcite[413--414 \& nn.45--46]{michelman1998}

Michelman 以 Michael H. v. Gerald D.\ix{cases}{02@美國!01@聯邦最高法院!1989Michael H. v. Gerald D.（1989）@Michael H. v. Gerald D.（1989）}（491 U.S. 110, 1989）為道德解讀的例證。案中未婚的生父請求探視子女，該子女出生時其母與另一男子有婚姻關係；州法長期否認生父在此情形下的探視權，生父主張其探視利益屬第十四修正案「liberty」〔自由〕所保護的內容。Scalia 的多數意見與 Brennan 的不同意見對照如下。\footcite[414]{michelman1998}

\begin{longtable}{@{}L{0.14\linewidth}L{0.38\linewidth}L{0.38\linewidth}@{}}
\caption{Michael H. v. Gerald D.：Scalia 多數意見與 Brennan 不同意見之方法對照。}\\
\toprule
項目 & Scalia（多數意見） & Brennan（不同意見） \\
\midrule
\endfirsthead
\toprule
項目 & Scalia（多數意見） & Brennan（不同意見） \\
\midrule
\endhead
結論 & 該探視利益不能是憲法所保護的利益 & 應檢視探視利益是否落入自由的一般原則 \\
依據 & 美國法律通常且傳統上從未給予處於聲請人地位的男子探視權（491 U.S. at 124） & 先前判決已把第十四修正案的「liberty」擴及社會傳統上視為重要的「generalized interests」〔概括性利益〕，包括「freedom from physical restraint, marriage, childbearing, [and] childrearing」〔免受身體拘束、婚姻、生育與養育子女〕（at 139） \\
\bottomrule
\end{longtable}

Brennan 認為決定性的問題，是生父的探視利益是否屬這些受傳統重視的概括性利益所例示的自由一般原則；法律向來拒絕回應該利益，可以解釋為法律未能符合其內在的理性標準，矯正這類失誤是他所任職位的主要任務之一。Michelman 評這個案件：「Nothing could better exemplify」〔再好不過地例示〕Dworkin 所推崇的道德解讀。\footcite[414]{michelman1998}

\section{合理多元與解釋之決定}
\label{sec:30_dworkin-4}

Michelman 在 414--415 頁接著指出，如 Dworkin 自己承認，道德解讀\ix{terms}{08@\zhlidxcat{憲法解釋與司法自制}!道德解讀}必然使解釋者介入無數涉及重大政治道德爭議的問題，Dworkin 稱之為「great and defining issues」〔重大而具決定性的議題〕，以及「intractable, profound questions of political morality that philosophers, statesmen, and citizens have debated for many centuries with no prospect of agreement」〔哲學家、政治家與公民辯論數世紀而無共識希望的棘手而深刻的政治道德問題〕（註 50、51，分別引 \en{Freedom's Law} 46 頁與 \en{Life's Dominion} 120 頁）。Michelman 的例子是他前文討論過的平等保護、正當程序與言論自由解釋：Adarand Constructors, Inc. v. Pena\ix{cases}{02@美國!01@聯邦最高法院!1995Adarand Constructors, Inc. v. Pena（1995）@Adarand Constructors, Inc. v. Pena（1995）}（515 U.S. 200, 1995）、Miller v. Johnson\ix{cases}{02@美國!01@聯邦最高法院!1995Miller v. Johnson（1995）@Miller v. Johnson（1995）}（515 U.S. 900, 1995）、Griswold v. Connecticut\ix{cases}{02@美國!01@聯邦最高法院!1965Griswold v. Connecticut（1965）@Griswold v. Connecticut（1965）}（381 U.S. 479, 1965）、Roe v. Wade\ix{cases}{02@美國!01@聯邦最高法院!1973Roe v. Wade（1973）@Roe v. Wade（1973）}（410 U.S. 113, 1973）、Buckley v. Valeo\ix{cases}{02@美國!01@聯邦最高法院!1976Buckley v. Valeo（1976）@Buckley v. Valeo（1976）}（424 U.S. 1, 1976）。\footcite[415 \& nn.3, 50--51]{michelman1998}

平等、民主、自由與公平的觀念可以幫助界定這類爭論的框架，最誠摯的承諾卻無從使答案不可置疑。Michelman 以 Rawls\ix{persons}{rawls@Rawls, John} 所說的合理多元\ix{terms}{02@\zhlidxcat{正當性與合理歧見}!合理多元}之「fact」〔事實〕與判斷之負擔\ix{terms}{02@\zhlidxcat{正當性與合理歧見}!判斷的負擔}（burdens of judgment）說明原因：人的經驗與見解各異，人的解釋判斷又受無數偶然因素左右（註 52，引 Rawls \en{Political Liberalism} 36--37、54--58 頁）。合理的歧見既然無從消除，這些解釋問題終須有人決定；「Whoever does not participate in the decision, for whatever reason, is to that extent governed by those who do」〔不論基於何種理由而未參與決定者，在此範圍內受參與決定者所治理〕。\footcite[415 \& n.52]{michelman1998}

由此，Michelman 說明道德解讀理論何以使其信奉者不能採程序性的辯護，亦即不能以人民制定法律的歷史事實證立憲法的民主性。依該理論，人民制定的內容在憲法操作意義之總和中所占份額太小，Brennan 這類法官對它的道德解讀所占份額太大；文本頒布與批准之時，留在政治上未決的意義太多。依自治觀，民主指人民以政治程序自行決定其有道德或物質理由關心的、政治上可決定的社會生活條件。因此，若有人要以道德解讀決定平權措施、種族歧視言論、政治獻金、任期限制、醫助自殺、墮胎、同志權利與槍枝管制這些問題，這個人應當是以民主方式行動的人民，「The main work of constitutional interpretation, cannot lie beyond democracy's purview」〔憲法解釋的主要工作，不能置於民主的管轄範圍之外〕（註 53 附於此句）。\footcite[415]{michelman1998}

\section{註 53 之三種反對}
\label{sec:30_dworkin-5}

註 53 自 415 頁頁底起，續於 416 頁頁底，記載 Michelman 對這個結論預期的抵抗。他預期抵抗「not much from Professor Dworkin」〔不多來自 Dworkin 教授〕，仍簡述如何應對最可能的三條反對路線。\footcite[415--416 n.53]{michelman1998}416 頁正文接著說：「Nevertheless it may have to, for the story is not over yet」〔然而它仍可能得如此，因為故事尚未結束〕。\footcite[416]{michelman1998}三項反對與回應排列如下。

\begin{longtable}{@{}L{0.14\linewidth}L{0.40\linewidth}L{0.38\linewidth}@{}}
\caption{Michelman 註 53 之三種反對及其回應。}\\
\toprule
反對 & 反對之內容 & Michelman 之回應 \\
\midrule
\endfirsthead
\toprule
反對 & 反對之內容 & Michelman 之回應 \\
\midrule
\endhead
abstraction（抽象） & 解釋者不難查證人民是否曾把某項抽象的政治道德原則（例如政治平等）制定入憲。若有，人民是否自治，取決於他們從法官等人收回的政府是否依該原則運作；只要依之，理論上人民即在自治 & 「I don't believe it」〔我不相信〕。該原則留待辯論與決定的問題，範圍與嚴重性足以嚴重損及自治；政府「colorblindness」〔色盲〕是否為不容妥協的憲法原則，是美國當時的爭論，這類爭議牽涉過多道德與物質利害，不能因人民自己寫下平等保護或正當程序條款而說憲法在此意義上具民主性（參原文第 III 部 C 節） \\
concretion（具體），或稱 integrity（整全性\ix{terms}{08@\zhlidxcat{憲法解釋與司法自制}!整全性}） & 解釋者的問題總是如何最佳綜合人民歷來制定入憲的多項抽象原則，再加上已證明持久的過往綜合性解釋。這項工作所受的資訊與拘束，比單一抽象原則的適用厚實 & 拘束仍不足以消除 Michelman 所提的問題。Dworkin 自己把這類綜合性的規範判斷描述為無底地政治性，且含美學判斷的成分（註 53 引 \en{Law's Empire} 73--76、87--93、229--32 頁）。誠摯而合理的爭論者對同一批材料必有多種建構，如此處見鴨、彼處見兔，或視為社會寫實悲劇、或視為嬉笑的影射小說（roman à clef），而人民不能把這類道德與物質利害重大的爭議交由他人決定，同時仍嚴肅主張自治 \\
right-answerism（唯一正解論\ix{terms}{08@\zhlidxcat{憲法解釋與司法自制}!抽象之反對／具體之反對／唯一正解論之反對}） & 憲法爭議在現實中無窮可爭，並不表示沒有答案正確與否之分。人民只要制定了原則的「grist」〔穀料〕，並由盡可能設計得當、足以產出真解釋的解釋「mill」〔磨坊〕處理，即屬自治 & 前提「without reservation」〔毫無保留〕接受；完整回應見原文第 V、VI 部（第 \ref{ch:40_paradox} 部） \\
\bottomrule
\end{longtable}

\section{Dworkin 論政治能動性}
\label{sec:30_dworkin-6}

\subsection{Michelman 所提之制度難題}

416 頁 C 節回到 Michelman 前文所稱的制度難題\ix{terms}{01@\zhlidxcat{民主理論}!制度難題}：個人如何能透過制度性地制定一項自己反對的法律而自治。他自問，為何不說對多數決立法的憲法限制支持積極自由\ix{terms}{01@\zhlidxcat{民主理論}!積極自由}，因為這類限制無損個人「control of his own fate」〔對自身命運的控制〕，在許多適用場合還可能增強之（註 54，引 \en{Freedom's Law} 21 頁）。\footcite[416 \& n.54]{michelman1998}Michelman 轉述 Dworkin 的回答：能誠實地說制定法律的主體只有一個，即人民或公民集體；個人要滿足自己作為自治之人的利益，唯一途徑是與該集體立法主體認同。制度難題的解決因此取決於一種可能性：個人見到周遭政治共同體以多數決政治程序立法，可藉把自己算作該共同體的一員，主張這些立法行為為己所有。\footcite[416--417 \& n.55]{michelman1998}

Michelman 在這處引了一段話，說民主是試圖藉由使政府決策服從於足以使公民產生參與、正當性\ix{terms}{02@\zhlidxcat{正當性與合理歧見}!正當性}與認同之感的溝通過程，調和個人自主與集體自決；即使個人與集體意志未必融合，公民仍可因參與這些溝通過程而把政府當作「their own」〔屬於自己的〕。接著他說這段話並非 Dworkin 所寫，作者是 Post\ix{persons}{post@Post, Robert}（註 56，引 \en{Constitutional Domains} 273 頁）。\footcite[417 \& n.56]{michelman1998}Post 的程序性民主理想及 Michelman 對它的批評，見第 2 部（\ref{ch:20_responsive}）。

\subsection{成員身分之三項條件}

Dworkin 為這種認同設定了合乎理性的前提：個人不能合理地把自己的政治能動性與任何集體結合，如果該集體的行動未對他道德與智識上的獨特性，以及他因而對集體結果的內容與對其影響集體結果之能力所懷的利益，維持應有的尊重。Michelman 轉述，集體須向每位成員保證下列三項（註 57，引 \en{Freedom's Law} 24--26 頁）：\footcite[417 \& n.57]{michelman1998}

\begin{enumerate}
\tightlist
\item 取得形成公共輿論的開放而有效管道，不受阻礙且平等；
\item 公共政策的決定中，每個人的利益受到同等程度的考慮；
\item 每個人自負責任地作道德與智識反思及判斷的能力，不受集體控制，受到隔絕保護。
\end{enumerate}

Michelman 認為這份清單看來很像美國憲法第一、第五與第十四修正案所保障者，確切地說，是 Brennan 這類憲法解釋者在這些條文中所發現者，而 Dworkin 的本意正是如此。Dworkin 的主張因此是：獨立司法機關正確解釋並實現憲法，可以確保個人把自己的政治能動性與政治共同體的立法行為認同時所需的合理前提，並藉此保障每個人在立法領域中自治的可能性，使違憲審查\ix{terms}{04@\zhlidxcat{違憲審查權限之歸屬}!違憲審查}若操作得宜，即解決制度難題\ix{terms}{01@\zhlidxcat{民主理論}!制度難題}。Michelman 稱之為「stunning argument」〔令人驚異的論證〕，並表示不認為它成功。\footcite[417]{michelman1998}

\subsection{belief 與 feeling 兩種讀法}

Michelman 的批評以認同須存在於個人意識為前提，意識表現為 belief（信念）或 feeling（感覺）。兩種讀法的檢討如下。\footcite[418]{michelman1998}

\begin{longtable}{@{}L{0.14\linewidth}L{0.40\linewidth}L{0.38\linewidth}@{}}
\caption{Michelman 對 Dworkin 認同說之兩種讀法的檢討（418 頁）。}\\
\toprule
讀法 & 內容 & Michelman 之批評 \\
\midrule
\endfirsthead
\toprule
讀法 & 內容 & Michelman 之批評 \\
\midrule
\endhead
belief & 若認同是一種信念，其內容不能是「我是立法者」或「我對立法結果有可察覺的影響」，因為 Dworkin 討論的起點之一，是大規模民主條件下沒有人能合理地這樣相信。信念的內容可以是個人有理由遵守集體所定之法 & 這項信念不足以滿足 Dworkin 的目的：它說明個人如何合理地尊重並接受他人所定之法，與說明個人如何合理地把自己視為自己的立法者，似乎相當不同 \\
feeling & 若認同是一種感覺，可以是個人對一個尊重其本人、獨立與利益的組織所作立法的滿足或驕傲，或是自己曾立法的感覺 & 這些感覺都不等於個人實際上立了法。依此讀法，Dworkin 似乎把 affect or attitude〔情感或態度〕誤當作 agency〔能動性〕。同情地認同行為人，不等於自己做了該行為 \\
\bottomrule
\end{longtable}

Michelman 承認，也許有人主張一種超過同情的、本體論意義的認同，使集體的行為得以「pass through」〔傳遞〕到成員身上，他指出這類有機體論的論調與 Dworkin 其餘的法律及政治哲學格格不入（註 58，引 \en{Taking Rights Seriously} 172 頁，該處把權利理論描述為以個人為中心）。\footcite[418 \& n.58]{michelman1998}

\subsection{Whitman 與 Brennan}

Michelman 引 Walt Whitman 自稱寬大而「contained multitudes」〔包容眾多〕，並說法官之中，Brennan 也寬大而包容眾多。包容眾多並不使 Brennan 成為每個人：「does not make Brennan everyone, or allow me to say that through his government I govern myself」〔不使 Brennan 成為每個人，也不使我得以說經由他的政府而自治〕。\footcite[418]{michelman1998}

\section{Dworkin 訊息之兩層}
\label{sec:30_dworkin-7}

Michelman 在 418 頁末說，Dworkin 的訊息表面令人愉快：它似乎使人能調和兩種最深的政治欲望，一是人人自治的民主願望，一是讓少數法官決定國家根本法的實際內容。Michelman 自己看到的是「deeper and darker-hued message」〔更深而色調更暗的訊息〕：沒有這種調和可得。接著回到前面提及的「crushing logical objection」〔壓倒性的邏輯反對〕，即不能把根本法的內容交由民主程序決定。\footcite[418--419]{michelman1998}

419 頁原文第 V 部開篇以此為出發點。Michelman 在該頁寫道：沒有任何程序能決定該程序自身是否適於決定這個問題或任何其他問題；並稱這是他所理解的 Dworkin 之「darker-hued message」：參與者須在平等、獨立、自由與安全的適當關係中進入程序，該程序才稱得上民主，因此民主程序之前必須先設定一份有操作內容、實質上等於權利清單\ix{terms}{03@\zhlidxcat{憲政主義與憲法之效力基礎}!權利法案}的規範（註 59 引 Dworkin：民主「cannot prescribe the procedures for testing whether the conditions for the procedures it does prescribe are met」〔不能規定檢驗其所規定之程序的條件是否具備的程序〕，\en{Freedom's Law} 33 頁）。\footcite[419 \& n.59]{michelman1998}這個悖論由 Michelman 第 V 部展開，見第 \ref{ch:40_paradox} 部。

\section{課前可提的問題}
\label{sec:30_dworkin-8}

\begin{enumerate}
\tightlist
\item Michelman 說 Dworkin 的內容檢驗「unassailable as far as it goes」，又說它對自治觀不足。對 accountability 觀的支持者而言，內容檢驗需要補上程序性的要件嗎？
\item 註 53 的 right-answerism 反對，Michelman 說前提「without reservation」接受，回應卻要到原文第 V、VI 部。他接受的前提，與他最終主張民主程序不能決定根本法內容，兩者如何並存？
\item Brennan 在 Michael H. 提出的「generalized interests」，與 Dworkin 所說從憲法傳統中提煉的政治道德原則，抽象程度是否相同？
\end{enumerate}
